\chapter{Хиральная кривизна Ходжа, ранговая неустойчивость и одна ядерная линия}

\section{Почему эта ветвь отличается от закрытой}

Чистая норма $\|\mathbb A^2\|^2$ не имела стационарного ненулевого фона.
Формульная генеалогия Томов IV--VI показывает другой механизм: квадрат
эрмитова отображения момента со сдвигом, который задан не после вычисления,
а градуировкой исходного комплекса. На физическом пакете одного поколения
имеется каноническое хиральное разложение
\begin{equation}
 H_{15}=H_L\oplus H_R,
 \qquad
 \dim_{\mathbb C}H_L=8,
 \qquad
 \dim_{\mathbb C}H_R=7.
 \label{eq:v7-h15-chiral-dimensions}
\end{equation}
Его разность размерностей равна единице и не использует класс $15$ как
подставленную цель.

Пусть $Y:H_L\to H_R$ является значением физической нечётной одноформы и
\begin{equation}
 d_Y=
 \begin{pmatrix}0&0\\Y&0\end{pmatrix},
 \qquad
 \Gamma_{15}=
 \begin{pmatrix}-I_8&0\\0&I_7\end{pmatrix}.
\end{equation}
Тогда канонический хиральный коммутатор равен
\begin{equation}
 [d_Y,d_Y^\dagger]
 =\begin{pmatrix}-Y^\dagger Y&0\\0&YY^\dagger\end{pmatrix}.
\end{equation}

\section{Новый заранее зарегистрированный функционал}

Определим
\begin{equation}
 \mathcal S_{\rm ch}(Y)
 =\frac1{15}\Tr_{H_{15}}
 \left([d_Y,d_Y^\dagger]-\Gamma_{15}\right)^2.
 \label{eq:v7-chiral-hodge-parent}
\end{equation}
Это не вычитание кривизны, выбранной по желаемому вакууму.
$\Gamma_{15}$ уже является частью конечной спектральной типизации и не
содержит юкавских коэффициентов, оси Хиггса или проектор будущего ядра.

Прямая блочная запись даёт
\begin{equation}
 \mathcal S_{\rm ch}(Y)
 =\frac1{15}\left\{
 \Tr_8(I_8-Y^\dagger Y)^2
 +\Tr_7(YY^\dagger-I_7)^2
 \right\}.
 \label{eq:v7-chiral-hodge-block-action}
\end{equation}
Если $\sigma_1,\ldots,\sigma_7$ --- сингулярные числа $Y$, то
\begin{equation}
 \mathcal S_{\rm ch}(Y)
 =\frac1{15}\left[
 1+2\sum_{j=1}^7(1-\sigma_j^2)^2
 \right]\ge\frac1{15}.
 \label{eq:v7-chiral-singular-value-action}
\end{equation}

\begin{theorem}[Канонический ранговый минимум]
Минимумы функционала~\eqref{eq:v7-chiral-hodge-parent} в полном пространстве
$\operatorname{Hom}(H_L,H_R)$ являются в точности коизометриями
\begin{equation}
 YY^\dagger=I_7,
 \qquad
 Y^\dagger Y=P_7,
 \qquad
 \rank P_7=7.
 \label{eq:v7-chiral-coisometry-vacuum}
\end{equation}
В каждом минимуме
\begin{equation}
 \dim_{\mathbb C}\ker Y=1,
 \qquad
 \min\mathcal S_{\rm ch}=\frac1{15}.
 \label{eq:v7-chiral-kernel-one}
\end{equation}
Остаточная энергия является точной ценой неравенства хиральных
размерностей, а не ручным дефектом.
\end{theorem}

\section{Нулевой сектор действительно неустойчив}

В отличие от ненулевого фонового $D_F$ предыдущего гейта, $Y=0$ является
стационарной точкой. Разложение имеет вид
\begin{equation}
 \mathcal S_{\rm ch}(Y)
 =1-\frac4{15}\|Y\|_{\rm HS}^2
 +\frac2{15}\Tr_8(Y^\dagger Y)^2.
 \label{eq:v7-chiral-zero-expansion}
\end{equation}
Поэтому на полном вещественном касательном пространстве размерности
$2\cdot8\cdot7=112$
\begin{equation}
 \Hess_0\mathcal S_{\rm ch}
 =-\frac8{15}I_{112}.
 \label{eq:v7-chiral-negative-hessian}
\end{equation}
Нулевая конфигурация стационарна, имеет отрицательные физические направления,
а положительный член четвёртой степени ограничивает действие снизу.

На радиальном пути $Y=tY_\star$ через любую коизометрию
\begin{equation}
 \mathcal S_{\rm ch}(tY_\star)
 =\frac1{15}\left[1+14(1-t^2)^2\right].
 \label{eq:v7-chiral-radial-path}
\end{equation}
Таким образом, отрицательная мода и её насыщение принадлежат одному
функционалу.

\section{Достижимость разрешёнными рёбрами Стандартной модели}

Разложим
\begin{equation}
 H_L=Q_L^{(6)}\oplus L_L^{(2)},
 \qquad
 H_R=u_R^{(3)}\oplus d_R^{(3)}\oplus e_R^{(1)}.
\end{equation}
Выбор значения одного хиггсовского дублета даёт ортогональные направления
$H$ и $\widetilde H$. Ребро $u$ отображает три верхних слабых цветовых
состояния в $u_R$, ребро $d$ --- три ортогональных слабых цветовых состояния
в $d_R$, а ребро $e$ --- заряженное лептонное направление в $e_R$.
Следовательно, существует разрешённое значение трёх рёбер, для которого
\begin{equation}
 Y_\star Y_\star^\dagger=I_7,
 \qquad
 \ker Y_\star=\mathbb C\nu_L.
 \label{eq:v7-physical-edge-coisometry}
\end{equation}
Машинный аудит строит такой оператор явно как матрицу $7\times8$.

Это пока не доказывает наблюдаемую нейтринную физику. Утверждается только,
что отсутствующее правое нейтринное ребро совпадает с единственной
неустранимой ядерной линией минимального хирального функционала.

\section{Граница результата}

Получен первый кандидат Тома VII, который одновременно имеет стационарный
нулевой сектор, отрицательный гессиан и ограниченный снизу ненулевой минимум.
Но это ещё хиральное ядро, а не полный родитель поля $\Phi$.

Открыты три обязательные проверки:
\begin{enumerate}
 \item вывести отображение $\Phi\mapsto Y_\Phi$ из полного модуля
 $E_{\rm aff}\widehat\otimes\mathcal Y_{\rm phys}$ без выбора линейного
 функционала на аффинном множителе;
 \item включить семейную размерность четыре и доказать, что две относительные
 свободы рёбер $u,d,e$ либо фиксируются тем же функционалом, либо честно
 остаются модами вакуумного многообразия;
 \item выполнить gauge/Real/junk/BRST--BV факторирование полного гессиана и
 отделить физические отрицательные направления от орбит симметрии.
\end{enumerate}

\begin{nogocriterion}[Запрет преждевременного нейтринного чтения]
Равенство $\dim\ker Y_\star=1$ нельзя называть выводом массы, смешивания или
майорановской природы нейтрино. Оно является только индексным следствием
хирального пакета $8\to7$ и разрешённых заряженных рёбер.
\end{nogocriterion}

\begin{protocol}
\begin{enumerate}
 \item канонический сдвиг из $\Gamma_{15}$: \textbf{получен};
 \item стационарность нулевого сектора: \textbf{получена};
 \item вещественных отрицательных направлений полного $7\times8$-ядра:
 \textbf{$112$};
 \item квартетное насыщение: \textbf{получено из того же функционала};
 \item минимальный ранг: \textbf{$7$};
 \item комплексная размерность ядра: \textbf{$1$};
 \item минимальная нормированная энергия: \textbf{$1/15$};
 \item достижимость рёбрами $u,d,e$: \textbf{получена конструктивно};
 \item подъём на полный аффинно-семейный модуль: \textbf{открыт};
 \item физическая частица: \textbf{не заявлена}.
\end{enumerate}
\end{protocol}

Машинный сертификат сохраняется в
\path{s2t_v7_chiral_hodge_index_instability_gate_results.json}.